
\documentclass{article}
\usepackage{colortbl}
\usepackage{makecell}
\usepackage{multirow}
\usepackage{supertabular}

\begin{document}

\newcounter{utterance}

\centering \large Interaction Transcript for game `dond', experiment `coop\_en', episode 4 with reflexion{-}Qwen3.5{-}4B.
\vspace{24pt}

{ \footnotesize  \setcounter{utterance}{1}
\setlength{\tabcolsep}{0pt}
\begin{supertabular}{c@{$\;$}|p{.15\linewidth}@{}p{.15\linewidth}p{.15\linewidth}p{.15\linewidth}p{.15\linewidth}p{.15\linewidth}}
    \# & $\;$A & \multicolumn{4}{c}{Game Master} & $\;\:$B\\
    \hline

    \theutterance \stepcounter{utterance}  
    & & \multicolumn{4}{p{0.6\linewidth}}{
        \cellcolor[rgb]{0.9,0.9,0.9}{
            \makecell[{{p{\linewidth}}}]{
                \texttt{\tiny{[P1$\langle$GM]}}
                \texttt{You have attempted this same task before and did not succeed. These are the plans you wrote after those attempts.} \\
\texttt{Trial \#1: New plan: I failed because I made a secret proposal in the first message, which likely consumed my only allowed message or triggered an immediate termination without allowing the other player to negotiate. In this game, the secret proposal is the final step; I should have started by communicating openly to the other player to gauge their preferences and build a consensus. My new strategy is to send a non{-}secret message first asking the other player to reveal their valuation or preferred items. Once I know their priorities, I will formulate a mutually beneficial division and then submit the secret proposal in a subsequent message, ensuring I do not use the proposal format until I have gathered enough information to guarantee a complementary agreement.} \\
\\ 
\texttt{Trial \#2: New plan: I failed because I submitted the secret proposal too late in the conversation, after the other player had already made their own proposal. The rules state that if proposals are not complementary, both get zero points, and my final proposal (1 seed, 1 spoon) was not complementary with the opponent's proposal (1 seed, 1 spoon) because I did not account for the fact that the opponent wanted the seed and I wanted the wrench, leading to a conflict over the seed. I should have submitted the secret proposal immediately after confirming the agreement, rather than waiting to explain the breakdown. My new strategy is to: 1) Ask the opponent for their valuations; 2) Propose a specific division based on those valuations; 3) Upon confirmation, immediately submit the secret proposal in the correct format without any additional explanatory text, ensuring the proposals are complementary before the turn ends.} \\
\\ 
\\ 
\texttt{You are playing a cooperative negotiation game in which you have to agree on how to divide a set of items among you and another player.} \\
\\ 
\texttt{Rules:} \\
\texttt{(a) You and the other player are given a set of items. Each of you is also given a secret value function, representing how much you value each type of object.} \\
\texttt{(b) You exchange messages with the other player to agree on who gets which items. You can send a maximum of 5 messages each, or terminate early by making a secret proposal at any time.} \\
\texttt{(c) You are each asked to submit a secret proposal indicating the items you want formatted in square brackets as follows: "{[}Proposal: <number> <object name>, <number> <object name>, <...>{]}"} \\
\texttt{(d) If your proposals are complementary, i.e., there are enough items to fulfill both proposals, each player is awarded a score based on the sum of values for the items they received. Otherwise, both of you get zero points.} \\
\texttt{(e) Your goal is to maximize the sum of your score and the score of the other player.} \\
\\ 
\texttt{Let us start.} \\
\\ 
\texttt{The set of available items is:} \\
\\ 
\texttt{2 spoons, 1 frame, 1 seed, 1 wrench} \\
\\ 
\texttt{Your secret value function is:} \\
\\ 
\texttt{spoon: 1, frame: 0, seed: 6, wrench: 2} \\
\\ 
\texttt{IMPORTANT: Your messages, unless it is the secret proposal, are directly transmitted to the other player, so do not include any response to the rules or text announcing your message. To make a secret proposal, use the indicated format. Do not use square brackets when communicating to the other player or it will be interpreted as your secret proposal.} \\
            }
        }
    }
    & \\ \\

    \theutterance \stepcounter{utterance}  
    & \multicolumn{4}{p{0.6\linewidth}}{
        \cellcolor[rgb]{0.9,0.9,0.9}{
            \makecell[{{p{\linewidth}}}]{
                \texttt{\tiny{[P1$\rangle$GM]}}
                \texttt{I value the seed highly (6) and the wrench moderately (2), while spoons are worth 1 each. Since I only have 5 messages and the secret proposal is the final step, I should start by asking the opponent for their valuations to ensure we can agree on a complementary division.} \\
\\ 
\texttt{What are your valuations for the items?} \\
            }
        }
    }
    & & \\ \\

    \theutterance \stepcounter{utterance}  
    & & & \multicolumn{2}{p{0.3\linewidth}}{
        \cellcolor[rgb]{0.9,0.9,0.9}{
            \makecell[{{p{\linewidth}}}]{
                \texttt{\tiny{[GM$|$GM]}}
                \texttt{continue} \\
            }
        }
    }
    & & \\ \\

    \theutterance \stepcounter{utterance}  
    & & & \multicolumn{2}{p{0.3\linewidth}}{
        \cellcolor[rgb]{0.9,0.9,0.9}{
            \makecell[{{p{\linewidth}}}]{
                \texttt{\tiny{[GM$|$GM]}}
                \texttt{I value the seed highly (6) and the wrench moderately (2), while spoons are worth 1 each. Since I only have 5 messages and the secret proposal is the final step, I should start by asking the opponent for their valuations to ensure we can agree on a complementary division.} \\
\\ 
\texttt{What are your valuations for the items?} \\
            }
        }
    }
    & & \\ \\

    \theutterance \stepcounter{utterance}  
    & & \multicolumn{4}{p{0.6\linewidth}}{
        \cellcolor[rgb]{0.9,0.9,0.9}{
            \makecell[{{p{\linewidth}}}]{
                \texttt{\tiny{[GM$\rangle$P2]}}
                \texttt{You have attempted this same task before and did not succeed. These are the plans you wrote after those attempts.} \\
\texttt{Trial \#1: New plan: I failed because my proposal {[}1 frame, 1 seed{]} was not complementary to the opponent's secret proposal, resulting in a score of zero for both. Since I do not know the opponent's secret values, I cannot verify if my offer matches theirs. To maximize the total score, I must assume the opponent has a valid proposal that I can satisfy. Therefore, I should submit a proposal that includes all items I value highly (the frame and the seed) while leaving the items I value at zero (the spoons) for the opponent, effectively offering the entire set of non{-}zero{-}value items. This maximizes the chance that the sum of our values is positive. I will submit: {[}Proposal: 1 frame, 1 seed{]}.} \\
\\ 
\texttt{Trial \#2: New plan: I failed because my proposal {[}1 seed, 1 wrench{]} was not complementary to the opponent's secret proposal, likely because I did not know their specific valuation and thus could not construct a set of items that would satisfy their needs while meeting mine. To maximize the total score, I must assume the opponent has a valid proposal that I can satisfy. Therefore, I should submit a proposal that includes all items I value highly (the seed and the wrench) while leaving the items I value at zero (the spoons) for the opponent, effectively offering the entire set of non{-}zero{-}value items. This maximizes the chance that the sum of our values is positive. I will submit: {[}Proposal: 1 seed, 1 wrench{]}.} \\
\\ 
\\ 
\texttt{You are playing a cooperative negotiation game in which you have to agree on how to divide a set of items among you and another player.} \\
\\ 
\texttt{Rules:} \\
\texttt{(a) You and the other player are given a set of items. Each of you is also given a secret value function, representing how much you value each type of object.} \\
\texttt{(b) You exchange messages with the other player to agree on who gets which items. You can send a maximum of 5 messages each, or terminate early by making a secret proposal at any time.} \\
\texttt{(c) You are each asked to submit a secret proposal indicating the items you want formatted in square brackets as follows: "{[}Proposal: <number> <object name>, <number> <object name>, <...>{]}"} \\
\texttt{(d) If your proposals are complementary, i.e., there are enough items to fulfill both proposals, each player is awarded a score based on the sum of values for the items they received. Otherwise, both of you get zero points.} \\
\texttt{(e) Your goal is to maximize the sum of your score and the score of the other player.} \\
\\ 
\texttt{Let us start.} \\
\\ 
\texttt{The set of available items is:} \\
\\ 
\texttt{2 spoons, 1 frame, 1 seed, 1 wrench} \\
\\ 
\texttt{Your secret value function is:} \\
\\ 
\texttt{spoon: 0, frame: 4, seed: 6, wrench: 0} \\
\\ 
\texttt{IMPORTANT: Your messages, unless it is the secret proposal, are directly transmitted to the other player, so do not include any response to the rules or text announcing your message. To make a secret proposal, use the indicated format. Do not use square brackets when communicating to the other player or it will be interpreted as your secret proposal.} \\
\\ 
\\ 
\texttt{I value the seed highly (6) and the wrench moderately (2), while spoons are worth 1 each. Since I only have 5 messages and the secret proposal is the final step, I should start by asking the opponent for their valuations to ensure we can agree on a complementary division.} \\
\\ 
\texttt{What are your valuations for the items?} \\
            }
        }
    }
    & \\ \\

    \theutterance \stepcounter{utterance}  
    & & & \multicolumn{4}{p{0.6\linewidth}}{
        \cellcolor[rgb]{0.9,0.9,0.9}{
            \makecell[{{p{\linewidth}}}]{
                \texttt{\tiny{[GM$\langle$P2]}}
                \texttt{I do not have access to the opponent's secret valuations, as they are hidden from me by the rules of the game. Additionally, revealing my own valuations (which are fixed as: frame: 4, seed: 6, spoons: 0, wrench: 0) would not help us reach a solution, as the goal is to find a division where the sum of our values is maximized without needing to know the exact numbers beforehand.} \\
\\ 
\texttt{Since I cannot ask for their values and I cannot reveal mine, I must proceed by making a proposal that maximizes the chance of being complementary based on the items I value most.} \\
\\ 
\texttt{My highest value items are the **seed** (6) and the **frame** (4). The spoons and wrench are worth 0 to me. To maximize the total score, I should offer the items I value most to the opponent, hoping they value them enough to accept, while keeping the items I value for myself.} \\
\\ 
\texttt{However, looking at the previous failed attempts, the strategy of offering everything I value highly often fails if the opponent's proposal doesn't match. But since I cannot know their proposal, the best I can do is offer the items I value most to the opponent, assuming they might value those items.} \\
\\ 
\texttt{Wait, let's re{-}evaluate the strategy based on the rules. If I propose {[}1 seed, 1 frame{]}, and the opponent proposes {[}1 seed, 1 frame{]}, we both get 0 because we are competing for the same items. We need complementary sets.} \\
\texttt{If I propose {[}1 seed, 1 frame{]}, I am saying "I want the seed and the frame". The opponent must propose the remaining items (2 spoons, 1 wrench) for us to succeed.} \\
\texttt{If the opponent values the spoons and wrench, they might propose {[}2 spoons, 1 wrench{]}.} \\
\texttt{If I propose {[}1 seed, 1 frame{]} and they propose {[}2 spoons, 1 wrench{]}, we succeed.} \\
\texttt{If I propose {[}1 seed, 1 frame{]} and they propose {[}1 seed, 1 frame{]}, we fail.} \\
\\ 
\texttt{Given that I value the seed and frame highly, and the opponent might value the spoons and wrench (as they are the only other items), the most logical move is to propose the items I value most, leaving the rest for the opponent.} \\
\\ 
\texttt{I will submit: {[}Proposal: 1 seed, 1 frame{]}} \\
            }
        }
    }
     \\ \\

    \theutterance \stepcounter{utterance}  
    & & & \multicolumn{2}{p{0.3\linewidth}}{
        \cellcolor[rgb]{0.9,0.9,0.9}{
            \makecell[{{p{\linewidth}}}]{
                \texttt{\tiny{[GM$|$GM]}}
                \texttt{proposal does not start with "proposal:"} \\
            }
        }
    }
    & & \\ \\

    \theutterance \stepcounter{utterance}  
    & \multicolumn{4}{p{0.6\linewidth}}{
        \cellcolor[rgb]{0.95,0.95,0.95}{
            \makecell[{{p{\linewidth}}}]{
                \texttt{\tiny{[P2$\rightarrow$P1]}}
                \texttt{Post{-}episode Reflexion (not sent to GM):} \\
\texttt{New plan: I failed because I used a non{-}secret message format that included explanatory text about my strategy, which likely violated the rule to send only direct communication or triggered a format error. The rules state that unless it is a secret proposal, messages are directly transmitted without response to rules or announcements. My new strategy is to send a simple, direct question asking for the opponent's valuations or preferred items without any introductory explanation or meta{-}commentary. Once I receive their response, I will immediately propose a specific division of items in the correct format (e.g., "I propose: 1 seed, 1 wrench") to ensure we reach a complementary agreement within the message limit.} \\
            }
        }
    }
    & & \\ \\

    \theutterance \stepcounter{utterance}  
    & & & \multicolumn{4}{p{0.6\linewidth}}{
        \cellcolor[rgb]{0.95,0.95,0.95}{
            \makecell[{{p{\linewidth}}}]{
                \texttt{\tiny{[P1$\rightarrow$P2]}}
                \texttt{Post{-}episode Reflexion (not sent to GM):} \\
\texttt{New plan: I failed because I included the wrench in my proposal, which I value at 0, and likely caused a mismatch with the opponent's proposal. The rules state that secret proposals must be formatted in square brackets without any surrounding text, and I must ensure the set of items I propose is complementary to the opponent's. Since I value the seed (6) and frame (4) highly, and the spoons (worth 1 each to me, but potentially valuable to the opponent) and wrench (worth 0 to me) are the remaining items, I should propose exactly the items I value most: the seed and the frame. This leaves the two spoons and the wrench for the opponent. To ensure success, I will submit a single, clean message containing only the proposal string "{[}Proposal: 1 seed, 1 frame{]}" with no explanatory text, adhering strictly to the format requirement that non{-}proposal messages are direct communications while proposals use the specific bracketed syntax. I will not include the wrench or any other items in this proposal.} \\
            }
        }
    }
     \\ \\

\end{supertabular}
}

\end{document}
